% Fragmento ES para inserción, no documento autónomo.
% Destino: V/main.tex, después de v_transporte_biparticion.tex y antes
% de 60_correlacion_cuasilibre_entropias.tex. ES, con pruebas contiguas.
\subsection{Énergie récupérable dans une distribution d’aire fixe}
\label{ins29:v:orientacion}

Le bord tritique permet de séparer le contenu de l’occupation totale
et celui de sa provenance sectorielle. Nous recevons les deux ensembles de
dix-huit liens, leurs occupations \(0,1,2\), les poids \(90/120\)
et la section thermique construite précédemment. Nous écrivons
\[
 X=N_{90},\qquad Y=N_{120},\qquad N=X+Y,\quad M=X-Y,
 \qquad d=30\Delta_{\rm mod}>0.
\]
Soit \(\epsilon=k_BT>0\) l’échelle énergétique de la réalisation
déclarée. Son Hamiltonien et son état de référence sont
\[
 H=\epsilon d(3X+4Y),\qquad
 \rho_0=Z^{-1}e^{-H/\epsilon}.
\]
L’opérateur d’aire déjà construit est
\(\widehat{\mathcal A}=a_{\rm b}N\), avec
\(a_{\rm b}=\gamma a_0L_\alpha^2>0\) fixe pendant le protocole.
Nous définissons \(W\) comme l’échange des deux ensembles de liens,
en conservant chaque occupation. Cette permutation inverse \(M\) et laisse
\(N\) invariant. L’inversion complémentaire des occupations traitée
précédemment, \(n\mapsto2-n\), conserve sa fonction propre et est une
opération distincte.

\begin{proposition}[Réponse énergétique invisible pour la lecture surfacique]
\label{ins29:v:respuesta}
L’état \(\rho_1=W\rho_0W^*\) a la même distribution complète
d’aire et la même entropie de von Neumann que \(\rho_0\), tandis que
\[
 \Delta E=\operatorname{Tr}[H(\rho_1-\rho_0)]
 =\epsilon\mathcal J(d)>0,
\]
\[
 \mathcal J(d)=D(\rho_1\Vert\rho_0)
 =18d\,[\mu(3d)-\mu(4d)],\qquad
 \mu(a)=\frac{e^{-a}+2e^{-2a}}{1+e^{-a}+e^{-2a}}.
\]
L’échange inverse récupère exactement \(\Delta E\) ; cette quantité
est l’énergie maximale récupérable au moyen d’une transformation unitaire
de l’état \(\rho_1\), pour le Hamiltonien fixe \(H\).
\end{proposition}
\begin{proof}
Comme \(WNW^*=N\), tous les projecteurs spectraux de l’aire commutent
avec \(W\) et conservent leurs probabilités. L’unitarité conserve le
spectre de l’état et son entropie. Dans la base des occupations, écrivons
\(p_0,p_1\) pour les probabilités. L’échange donne
\[
 \log\frac{p_1}{p_0}=-dM,\qquad
 \langle M\rangle_{p_1}=-\langle M\rangle_{p_0},\qquad
 \langle M\rangle_{p_0}=18[\mu(3d)-\mu(4d)].
\]
La somme pondérée du logarithme prouve la formule de
\(D(\rho_1\Vert\rho_0)\). En outre,
\(\mu'(a)=-\operatorname{Var}_a(n)<0\) pour les trois poids positifs,
de sorte que \(\mathcal J(d)>0\).
L’identité de Gibbs, avec l’égalité des entropies, donne
\(\Delta E=\epsilon D(\rho_1\Vert\rho_0)\).
Pour tout unitaire \(V\), l’état \(V\rho_1V^*\) conserve cette
même entropie et satisfait
\[
 \operatorname{Tr}[H(V\rho_1V^*-\rho_0)]
 =\epsilon D(V\rho_1V^*\Vert\rho_0)\geq0.
\]
Le choix \(V=W^*\) atteint l’égalité et prouve le maximum.
\end{proof}

La distribution surfacique conserve toutes les multiplicités de \(N\),
mais le Hamiltonien distingue leurs répartitions \((X,Y)\). L’excès
énergétique se situe exactement dans cette distinction. L’échelle
d’horizon, lorsque la réalisation correspondante est utilisée, est
\(\epsilon=\hbar c/(2\pi R)\). Ainsi, l’action transforme
\(\mathcal J(d)\) en une réponse énergétique à rayon fixe.
L’opérateur d’aire de cette préparation et l’aire géométrique d’un horizon
avec rétroaction conservent leurs domaines respectifs.

\subsubsection{Restitution conditionnée et fonctionnelle réduite}

La localisation intrafibre admet une formulation exacte. Sur
l’espace fini des configurations, soit \(q\) un extracteur surjectif,
\(p_{0,q}=q_*p_0\) et \(p_q=q_*p\). Pour une distribution \(r\) de
l’extracteur, nous définissons
\[
 (\mathcal R_0r)(x)
 =r(q(x))\frac{p_0(x)}{p_{0,q}(q(x))}.
\]
Ce restituteur satisfait \(q_*\mathcal R_0=I\). Sa distribution
conditionnelle conserve les poids thermiques à l’intérieur de chaque fibre.

\begin{proposition}[Énergie libre de l’information intrafibre]
\label{ins29:v:restitucion}
Pour \(F_T(p)=\sum_xp(x)H(x)+\epsilon\sum_xp(x)\log p(x)\),
\[
 F_T(p)-F_T(\mathcal R_0p_q)
 =\epsilon D(p\Vert\mathcal R_0p_q).
\]
Par conséquent, \(\mathcal R_0p_q\) est l’unique minimiseur de l’énergie
libre parmi les distributions de marginale \(p_q\).
\end{proposition}
\begin{proof}
Sur le support de \(p\), la factorisation
\[
 \frac{p(x)}{p_0(x)}=
 \frac{p_q(q(x))}{p_{0,q}(q(x))}
 \frac{p(x)}{(\mathcal R_0p_q)(x)}
\]
donne, en prenant les logarithmes et en sommant,
\(D(p\Vert p_0)=D(p_q\Vert p_{0,q})+
D(p\Vert\mathcal R_0p_q)\).
L’identité \(F_T(p)=\epsilon D(p\Vert p_0)-\epsilon\log Z\)
et la même formule appliquée à \(\mathcal R_0p_q\) prouvent
l’énoncé. La divergence s’annule exactement lorsque ses arguments
coïncident ; on adopte \(0\log0=0\).
\end{proof}

Pour \(q=N\), écrivons
\(g_{18}(x)=[t^x](1+t+t^2)^{18}\), la multiplicité tritique déjà
construite. Le Hamiltonien effectif de la lecture surfacique est
\[
 \frac{H_N(n)}\epsilon=
 \frac{7d}{2}n-\log\!\sum_{x+y=n}
 g_{18}(x)g_{18}(y)e^{d(x-y)/2}.
\]
Cette expression provient de la somme de \(e^{-H/\epsilon}\) sur la fibre.
La somme est invariante sous l’inversion du signe du terme orienté,
par échange de \(x\) et \(y\). Pour la préparation précédente,
\(p_{1,N}=p_{0,N}\) et \(\mathcal R_0p_{1,N}=p_0\) : tout l’excès
d’énergie libre réside dans l’information que cette lecture surfacique omet.
En revanche, conserver \((N,M)\) fixe l’énergie de chaque fibre et rend
uniforme sa distribution conditionnelle. Le résultat précise quels
observables suffisent pour récupérer cette collection et sa réponse thermique.
